\section{Model Assumptions}

To simplify the problem while preserving physical accuracy, we make the following assumptions:

\begin{enumerate}
  \item \textbf{Smooth planar interfaces:} All film interfaces are assumed to be perfectly smooth and parallel. Surface roughness effects on optical scattering are neglected, as RMS roughness is typically much smaller than the relevant wavelengths.

  \item \textbf{Isotropic and homogeneous media:} Each layer material is assumed to be optically isotropic and homogeneous. The complex refractive index $\tilde{n}(\lambda) = n(\lambda) + i k(\lambda)$ fully characterizes the optical response at each wavelength.

  \item \textbf{Coherent thin-film interference:} For film thicknesses below the coherence length ($\lesssim 100$ $\mu$m), optical interference effects are treated coherently using the Airy summation formula. For very thick layers, the incoherent limit is approached.

  \item \textbf{Normal and near-normal incidence approximation:} For radiative transfer calculations, we integrate over the hemisphere using Gauss-Legendre quadrature to account for angular dependence of emissivity.

  \item \textbf{Kirchhoff's law applicability:} We invoke Kirchhoff's law of thermal radiation, equating spectral directional absorptivity to spectral directional emissivity: $\varepsilon(\lambda, \theta) = A(\lambda, \theta) = 1 - R(\lambda, \theta) - T(\lambda, \theta)$.

  \item \textbf{Steady-state thermal conditions:} The radiative cooler is assumed to operate under steady-state thermal equilibrium, with net cooling power $P_{\text{cool}} = 0$ determining the equilibrium temperature.

  \item \textbf{Standard atmospheric and solar conditions:} Atmospheric transmittance follows the MODTRAN mid-latitude summer model. Solar irradiance follows the ASTM G173-03 AM1.5 Global Tilt spectrum (1000 W/m$^2$ total).

  \item \textbf{Non-radiative heat transfer:} Conduction and convection are modeled as a lumped linear term $P_{\text{nonrad}} = h_c(T_{\text{amb}} - T)$, with $h_c = 6.9$ W/(m$^2\cdot$K) for natural convection.
\end{enumerate}
